\documentclass{article}

\usepackage[utf8]{inputenc}

\usepackage{amsmath, bm}
\usepackage{graphicx}
\usepackage{amssymb}
\usepackage{float}
\usepackage{caption}
\usepackage{subcaption}
\usepackage{hyperref}
\usepackage{tikz}
\usepackage{layout}

\usepackage[margin=1in]{geometry}
\usepackage{listings}
\usepackage{xcolor}
\usepackage{color, colortbl}
\usepackage{textgreek}
\usepackage{mathrsfs}
\usepackage{savetrees}

\usetikzlibrary{calc}
\usetikzlibrary{angles,quotes} % for pic
\usetikzlibrary{patterns,snakes}
\usetikzlibrary{arrows}
\tikzset{>=latex} % for LaTeX arrow head

\setlength{\parskip}{\baselineskip}%
\setlength{\parindent}{0pt}%
\linespread{0.9}


\definecolor{codegreen}{rgb}{0,0.6,0}
\definecolor{codegray}{rgb}{0.5,0.5,0.5}
\definecolor{codepurple}{rgb}{0.58,0,0.82}
\definecolor{backcolour}{rgb}{0.95,0.95,0.92}

\lstdefinestyle{mystyle}{
    backgroundcolor=\color{backcolour},   
    commentstyle=\color{codegreen},
    keywordstyle=\color{magenta},
    numberstyle=\tiny\color{codegray},
    stringstyle=\color{codepurple},
    basicstyle=\ttfamily\footnotesize,
    breakatwhitespace=false,         
    breaklines=true,                 
    captionpos=b,                    
    keepspaces=true,                 
    numbers=left,                    
    numbersep=5pt,                  
    showspaces=false,                
    showstringspaces=false,
    showtabs=false,                  
    tabsize=2
}

\lstset{style=mystyle}



\begin{document}

\title{Acoustic Signatures of Quadcopter Propellers \\ Research Design and Methodology}
\author{lwp26}
\date{October 2024}
\maketitle 

\iffalse
\begin{abstract}
    \centering
    This report presents the development of a software tool to help solve constrained design optimization of a shell and tube heat exchanger.
    The software is developed in Python and uses LMTD and NTU methods to solve for the output parameters.
    Computational model is compared with experimental data 
    Various optimization algorithms were tested with constraints and the results are presented.
\end{abstract}
\fi

%-----------------------------------------------------------------------------------------
\section{Introduction}
%-----------------------------------------------------------------------------------------

\subsection{Purpose}
Drone based services are becoming increasingly popular in various industries such as agriculture and delivery services [?].
However the adoption of drone services has been limited by public perception and legislation [?].
A key factor influencing public perception is the noise generated by drones, which is often perceived as intrusive and annoying [?].
Propulsive efficiency is also important for feasibility as it directly affects flight time and payload capacity [?].
This project aims to investigate acoustic sources of noise in quadcopter propellers and reduce the noise generated by these propellers.


% design and theory of test rig and data acquisition system
% precise problem statement of experiment and how it fits into reducing noise


\subsection{Objectives}
% set goal of what to achieve

The main objective is to design a test rig to measure the acoustic signature of a single propeller design.

This however, can be broken down into smaller objectives:

\begin{itemize}
    \item Research and document required components
    \item Decide using available theory:
    \begin{itemize}
        \item how best to orient load cells to measure thrust and torque
        \item where best to place microphones to reconstruct the sound field
        \item how to process the data to extract useful information
    \end{itemize}
    \item Create CAD model of test rig
    \item Write calibration and testing procedures
\end{itemize}


% bibliography

\begin{thebibliography}{9}

%Endres, SC, Sandrock, C, Focke, WW (2018) “A simplicial homology algorithm for lipschitz optimisation”, Journal of Global Optimization.

  \bibitem{handout}
  J. V. Taylor and J. C. Massey
  \emph{GA3 Heat Exchanger Handout}
  University of Cambridge,
  2024.

\end{thebibliography}

\end{document}